% ============================================================================
% ABSTRACT
% BU guide 1.5.1: the abstract must contain a clear and brief statement of
%   (1) the problem, (2) the procedures / methods followed, (3) the results,
%   and (4) the conclusions.
% There is no word limit any more, but be concise. Double spaced.
% Do NOT include graphs, charts, tables or illustrations here (guide 1.5.3).
% Include full proper nouns and place names -- the abstract is keyword
% searchable in ProQuest.
% ============================================================================

Apical dendrites of layer 5 pyramidal cells (L5PCs) are critical integrative hubs in cortical circuits. 
The repertoire of their activity patterns, including $\mathrm{Ca}^{2+}$ and NMDA spikes, varies across brain states.
However, the relationship between neuromodulatory activity and the statistical properties of dendritic $\mathrm{Ca}^{2+}$ and NMDA spikes remains unclear.
In this study, we performed volumetric imaging to characterize spontaneous $\mathrm{Ca}^{2+}$ dynamics in the apical dendrites of L5PCs in awake, head-fixed mice. 
We used Rbp4-Cre mice, retro-orbitally injected with PHP.eB AAVs to express GCaMP6s in L5PCs, and a red fluorescent acetylcholine (ACh) sensor GRAB-ACh. 
Using light sheet microscopy (SCAPE-2.0, Voleti et al. 2019) we imaged spontaneous $\mathrm{Ca}^{2+}$ and ACh activity in the primary somatosensory (SI) cortex at a volumetric rate of $\SI{5}{\hertz}$ with a $400 \times 200 \times 150 \,\mathrm{\mu m}^3$ field of view, capturing the top $\approx 150 \,\mathrm{\mu m}$.
Simultaneously, we acquired behavior readouts: the eye pupil diameter, whisking, and accelerometer-based body movement. 
Data were analyzed using a custom Python toolbox for detection and segmentation of dendritic events from volumetric time series.
Our results show high correlation between global (averaged within the volume) dendritic $\mathrm{Ca}^{2+}$ dynamics, ACh, and the behavior readouts.
Activity-dependent segmentation of dendritic regions of interest (ROIs) revealed several types of events including vertically oriented putative $\mathrm{Ca}^{2+}$ spikes engulfing the nexus and apical trunk as well as $\mathrm{Ca}^{2+}$ transients localized to specific dendritic branches (putative NMDA-dependent synaptic events).
The probability of putative $\mathrm{Ca}^{2+}$ spikes increased with an increase in ACh. 
To investigate the relationship between dendritic $\mathrm{Ca}^{2+}$ activity and somatic spiking, we used random access 2-photon microscopy (Femto3D Atlas, Femtonics). 
Our preliminary data suggest that most of the events observable within the top $\approx 150 \,\mathrm{\mu m}$ correspond to synaptic inputs and dendritic $\mathrm{Ca}^{2+}$ spikes rather than back-propagated action potentials.
